% =========================================================
% SECCIÓN 2.2 - DESCRIPCIÓN DEL SISTEMA EXPERIMENTAL
% =========================================================

El diseño experimental se estructuró como un ensayo de laboratorio controlado en un reactor batch cerrados, orientado a cuantificar las emisiones de \ce{CO2} y \ce{CH4} durante el ciclo larval de \textit{Hermetia illucens} bajo condiciones ambientales estandarizadas. La elección de un sistema batch responde a la necesidad de mantener trazabilidad de masa (sustrato, biomasa, gases) en ventanas temporales discretas, facilitando la correlación posterior con el modelo bioenergético del Capítulo 3 \parencite{eriksenDynamicModellingFeed2022, padmanabhaComprehensiveDynamicGrowth2020}. A continuación se detallan las especificaciones del reactor, las condiciones de laboratorio y el manejo de la biomasa y el sustrato.

% ---------------------------------------------------------
\subsection{Reactor batch y condiciones de laboratorio}
\label{subsec:cap2_reactor}

Los ensayos se ejecutaron en un contenedor de polipropileno alimentario tipo \emph{tupper} o envase de conservación doméstica, con dimensiones aproximadas de $12\times12\times5$~cm, lo cual brinda una capacidad útil estimada de \SI{0,7}{\litre} a \SI{0,9}{\litre}. 
El sistema se operó como reactor batch cerrado mediante tapa de plástico con sellado por presión. 
Se insertaron sondas de sensores NDIR para \ce{CO2} (modelo Winsen MH-410D, rango 0--\SI{5000}{ppm}, resolución $\pm$\,\SI{50}{ppm}) y \ce{CH4} (modelo Winsen MH-441D, rango 0--\SI{10}{\percent} en volumen, resolución \SI{0,01}{\percent}), así como una sonda de temperatura/humedad relativa (\cref{fig:cap2_reactor_esquema}). 
El \emph{headspace} o volumen de cabeza libre ---esto es, el espacio de aire comprendido entre la superficie del sustrato y la tapa del contenedor--- se estandarizó en aproximadamente \SI{0,3}{\litre} a \SI{0,4}{\litre}, asumiendo una carga de sustrato de \SI{250}{\gram} que ocupa aproximadamente \SI{2}{\centi\metre} de altura, con el fin de mantener comparable la acumulación gaseosa durante las ventanas cerradas de medición. 
Como los tiempos de medición fueron cortos, se registró la temperatura del entorno y se constató que era aproximadamente constante, manteniendo un valor de \SI{30}{\celsius} $\pm$\,\SI{2}{\celsius} y una humedad relativa (HR) del \SI{60}{\percent}--\SI{80}{\percent}, rangos reportados como óptimos para la maximización de la tasa metabólica larval \parencite{salamEffectDifferentEnvironmental2022, nayakHermetiaIllucensDiptera2024}.

El protocolo de medición se estructuró en eventos discretos espaciados cada 2 días. 
En cada evento, se cerraba herméticamente el reactor y se registraba la acumulación gaseosa mediante los sensores NDIR hasta alcanzar la saturación de los instrumentos. 
Una vez saturados los sensores, se abría el reactor para restaurar las condiciones ambientales y continuar el ciclo de bioconversión. 
Este procedimiento se repitió 6--7 veces a lo largo de un periodo experimental de 12--\SI{15}{\day}, tras el cual las larvas alcanzaron el estadio de pre-pupa..

% ---------------------------------------------------------
\subsection{Biomasa larval y manejo del sustrato}
\label{subsec:cap2_biomassa_sustrato}

Las larvas de \emph{Hermetia illucens} se obtuvieron de una colonia estabilizada en la Sede Palmira (Universidad Nacional de Colombia), sincronizadas por edad (5~d post-eclosión, estadio L3 inicial) y pesadas en grupo para determinar la biomasa inicial húmeda. 
La densidad de siembra se fijó en aproximadamente 700 larvas por reactor, distribuidas sobre una carga de \SI{250}{\gram} de sustrato. 

El sustrato consistió en dos tratamientos experimentales: (i) residuos de procesamiento avícola (pollinaza) y (ii) una mezcla homogeneizada de residuos agroindustriales (pulpa de café, \emph{Coffea arabica}, y cáscaras de frutas tropicales, plátano y papaya, en proporción 70:30 base húmeda), seleccionada por su alta aptitud para la bioconversión y disponibilidad estacional en el Valle del Cauca \parencite{brunoValorizationOrganicWaste2025}. 
En ambos tratamientos, la masa de sustrato se ajustó a \SI{250}{\gram} por reactor. 
La humedad inicial del sustrato se reguló a $65\%\pm5\%$ mediante adición controlada de agua destilada, valor dentro del intervalo 60\%--80\% que \textcite{bekkerImpactSubstrateMoisture2021} identifican como crítico para el mantenimiento de la tasa de asimilación y el crecimiento exponencial larval. 
El pH se registró entre 5.5 y 6.5, sin corrección química, dado que la acidificación natural por metabolismo microbiano no alcanzó niveles inhibitorios durante el ciclo.

Se establecieron seis réplicas ($n = 6$) por condición experimental, más dos controles (sustrato sin larvas), distribuidos en bloques completos al azar para controlar efectos de posición. 
No se ejecutó volteo mecánico del sustrato, salvo en las réplicas destinadas al estudio de perturbación (Capítulo~5), para preservar la formación de gradientes naturales de oxigenación que permiten evaluar la validez del modelo DEB bajo heterogeneidad real. 
Al final del ciclo (12--\SI{15}{\day}), se procedió a la separación mecánica de larvas del sustrato residual (frass), pesaje individual de una submuestra ($n=30$ por réplica) y determinación de materia seca ($60^{\circ}$C, 48~h) para expresar la biomasa en base seca.